\section{Retaining the unused four-body blocks}
\label{app:retained-blocks}
\begingroup
\newcommand{\rbnorm}[1]{\lVert #1\rVert}
\newcommand{\rbket}[1]{\lvert #1\rangle}
\newcommand{\rbbra}[1]{\langle #1\rvert}
\newcommand{\rbip}[2]{\langle #1,#2\rangle}

The local loss construction requires a stronger operator inequality than
the unperturbed gap: most of the normal-ordered four-body term must remain
on the right-hand side. We give the full finite-dimensional calculation
that establishes this strengthening. The seven auxiliary squares, their
rational certificate, and the finite-flux transfer are adapted from
\cite[Sections 2--6 and 8]{Gap}. The proof below retains the
four-body blocks that are discarded when only the gap is sought.

We use the lift and wedge maps of \eqref{eq:lift} and
\eqref{eq:wedge-maps}, with increasing orthonormal wedges of norm one.
Thus $B(v\wedge w)=B(w)B(v)$, and the domain of $W_4$ uses ordered
pair factors. Write $a=q-1$ and $H=H_q$; the conventions give
\begin{align}
 \mathcal L_k(\rbket v\rbbra w)&=B(v)^\dagger B(w),
       \label{rb:eq:lift}\\
 B_p=B(v_p),\qquad H&=\mathcal L_2(\Pi_{V_q})
            =\sum_{p=0}^{2q-2}B_p^\dagger B_p.
       \label{rb:eq:fockH}
\end{align}
The dagger in this appendix denotes the same adjoint as the star in the
main text.

The tensor product $V_q\otimes V_q$ contains one copy of each spin
$2q-2-l$, $0\le l\le 2q-2$. Let $P_{q,l}$ be its orthogonal
spin projection and define
\begin{equation}\label{rb:eq:retained-projections}
 R_q^{\mathrm{lo}}=\sum_{l=0}^{22}P_{q,l},\qquad
 R_q^{\mathrm{hi}}=I-R_q^{\mathrm{lo}}\quad(q\ge24).
\end{equation}
The spin-block label used in the calculation below is $D=l+1$;
thus $R_q^{\mathrm{lo}}$ consists exactly of $1\le D\le23$.

\begin{proposition}[Retained four-body blocks]\label{prop:retained-blocks}
For every fixed $0<c_0<1/25$ there exists $q_0$ such that, for every integer
$q\ge q_0$, the following inequality holds on the entire Fock space
$\mathcal F_q$:
\begin{equation}\label{rb:eq:retained-blocks}
 H_q^2\ge c_0H_q+
 \mathcal L_4\!\left(W_4R_q^{\mathrm{hi}}W_4^\dagger\right).
\end{equation}
The threshold is independent of particle number.
\end{proposition}

We prove the proposition by comparing a positive rotational average of
seven explicit squares with the normal-ordered terms of $H_q^2$.
The comparison uses the target four-body blocks $1\le D\le23$ only.
All other target blocks consequently remain with coefficient one.

\begin{lemma}[Normal-ordered square]\label{rb:lem:square}
On $\mathcal F_q$,
\begin{equation}\label{rb:eq:square}
 H^2=H+\mathcal L_3\!\left(W_3(W_3^\dagger W_3-I)W_3^\dagger\right)
       +\mathcal L_4(W_4W_4^\dagger).
\end{equation}
\end{lemma}
\begin{proof}
In normal ordering $B_p^\dagger(B_pB_r^\dagger)B_r$, only the two inner annihilators cross the two inner creators. Two, one, or zero contractions leave two-, three-, or four-body terms, respectively. On two particles, $H=\Pi_{V_q}$ is a projection, so the two-body coefficient is $\Pi_{V_q}$. The term with no contractions is
\[
 \sum_{p,r}B(v_p\wedge v_r)^\dagger B(v_p\wedge v_r)
 =\mathcal L_4(W_4W_4^\dagger).
\]
There is no extra factor of two: this sum, and the defining basis of $V_q\otimes V_q$, both use ordered pairs. On three particles, $H=W_3W_3^\dagger$, since $W_3$ restricted to $v_p\otimes U_q$ is the map $B_p^\dagger$ from one to three particles. Subtracting its two-body contribution from $H^2$ therefore determines the remaining coefficient as
\[
 (W_3W_3^\dagger)^2-W_3W_3^\dagger
 =W_3(W_3^\dagger W_3-I)W_3^\dagger.
\]
Normal-ordered coefficients are determined successively by these restrictions, proving the identity on Fock space.
\end{proof}

\subsection{Spin coefficients and their limits}
\label{rb:sec:spin}
The three-body comparison will use the exact spin spectrum of the wedge map. The four-body comparison will use only finitely many coupling coefficients and their limits.

For this auxiliary calculation we simultaneously change the signs of all pair rows from \eqref{eq:pair-rows}, so that $v_0=-e_0\wedge e_1$, and use the normalized lowering bases $e_p$ and $v_p$. This common phase change leaves $H_q$ and the intrinsic wedge maps and spin projections unchanged. For spin $j$, lowering from label $p$ to $p+1$ has coefficient $\sqrt{(p+1)(2j-p)}$. Write $(s)_p=s(s-1)\cdots(s-p+1)$, with $(s)_0=1$.

\begin{lemma}[Pair coefficients]\label{rb:lem:pair}
The coefficient of $e_i\wedge e_j$ in $v_p$, for $i<j$, vanishes unless $i+j=p+1$, and otherwise equals
\begin{equation}\label{rb:eq:pair}
 v_p(i,j)=(i-j)\sqrt{\frac{(q)_i(q)_j\,p!}{q(2q-2)_p\,i!j!}}.
\end{equation}
For fixed $p,i,j$, its limit is
\begin{equation}\label{rb:eq:pairstar}
 v_p^*(i,j)=(i-j)\sqrt{\frac{p!}{2^p i!j!}}\,\mathbf1_{i+j=p+1}.
\end{equation}
\end{lemma}
\begin{proof}
Apply simultaneous lowering $p$ times to $-e_0\wedge e_1$ and divide by $\sqrt{(2q-2)_p p!}$. Expand the lowering operator binomially. Combining the two contributions to each increasing pair gives \eqref{rb:eq:pair}. Taking the limit gives \eqref{rb:eq:pairstar}.
The normalization can also be checked directly: the square of the
coefficient in \eqref{rb:eq:pair} is
\[
 (i-j)^2\frac{\binom qi\binom qj}{q\binom{2q-2}{p}},
\]
and differentiating \((1+x)^q\) gives
\[
 \sum_{i,j=0}^q(i-j)^2\binom qi\binom qj x^{i+j}
 =2qx(1+x)^{2q-2}.
\]
Take the coefficient of \(x^{p+1}\) and halve the sum to restrict to \(i<j\).
\end{proof}

The elementary tensor product rule decomposes spins $j_1,j_2$ into one copy of each spin $j_1+j_2-z$, for $0\le z\le\min(2j_1,2j_2)$. Let $R_{3,z}$ be the corresponding projection in $V_q\otimes U_q$, and put
\[
 d_a=2q-1,\qquad d_{3,z}=3q-1-2z.
\]
For $q\ge2$, the allowed indices are $0\le z\le q$.

\begin{lemma}[Three-body Gram eigenvalues]\label{rb:lem:q}
\begin{equation}\label{rb:eq:q}
 W_3^\dagger W_3-I=\sum_{z=0}^q q_z(q)R_{3,z},\qquad
 q_z(q)=(-1)^z\frac{(q)_z}{(2q-2)_z}
          \left(3z-1-\frac{z(z+1)}q\right).
\end{equation}
In particular $W_3R_{3,z}=0$ for $z=0,1,3$ whenever these indices are allowed. For fixed $z$,
\begin{equation}\label{rb:eq:qlimit}
 q_z(q)\longrightarrow (3z-1)(-1/2)^z.
\end{equation}
\end{lemma}
\begin{proof}
Identify $V_q\otimes U_q$ with tensors antisymmetric in their first two positions. Then $W_3$ is $\sqrt3$ times full antisymmetrization. Its Gram operator is the compression of $I-P_{13}-P_{23}$, where $P_{ij}$ swaps tensor positions. The two swap compressions coincide, since conjugation by $P_{12}$ interchanges them and $P_{12}=-I$ on the domain. By equivariance and multiplicity-freeness they act by the same scalar $b_z$ on each spin block.

A highest vector of that block has coefficients $s_p$ on $v_p\otimes e_{z-p}$, $0\le p\le z$, satisfying
\begin{equation}\label{rb:eq:highest-ratio}
 \frac{s_{p+1}}{s_p}=-\sqrt{\frac{(z-p)(q-z+p+1)}{(p+1)(2q-2-p)}}.
\end{equation}
This follows by applying the raising operator. Using \eqref{rb:eq:pair}, the coefficient of the compressed $P_{23}$ image on $v_0\otimes e_z$ is
\[
 s_0b_z=\frac12\left[
 s_{z-1}\sqrt{\frac{z(q)_z}{q(2q-2)_{z-1}}}
 -s_z(z-1)\sqrt{\frac{(q)_z}{(2q-2)_z}}\right],
\]
with the first term absent for $z=0$. Here the ordered tensor coefficients of a normalized wedge include a factor $1/\sqrt2$. Iteration of \eqref{rb:eq:highest-ratio} gives
\[
 \frac{s_z}{s_0}=(-1)^z\sqrt{\frac{(q)_z}{(2q-2)_z}},\qquad
 \frac{s_{z-1}}{s_0}=(-1)^{z-1}
 \sqrt{\frac{z(q-1)_{z-1}}{(2q-2)_{z-1}}}\quad(z\ge1).
\]
Substitution gives $q_z=-2b_z$ in the form \eqref{rb:eq:q}. The three stated null sectors have $q_z=-1$, and \eqref{rb:eq:qlimit} follows directly.
\end{proof}

We will repeatedly use a fixed-deficit limit of this same calculation. In a product of spins $j_1,j_2$, total lowering deficit means $T=j_1+j_2-m$, where $m$ is the total magnetic quantum number. Coefficients below are with respect to normalized monomials $X^pY^{T-p}/\sqrt{p!(T-p)!}$, and are zero outside the stated index ranges.

\begin{lemma}[Fixed-deficit coupling limit]\label{rb:lem:coupling}
Suppose $j_1,j_2\to\infty$ and $j_2/(j_1+j_2)\to u^2$, where $0<u<1$. Put $v=\sqrt{1-u^2}$. For fixed integers $0\le z\le T$, the normalized coupled vector of spin $j_1+j_2-z$ at total lowering deficit $T$ can be chosen with real coefficients converging to the coefficients $s^{(u)}_{z,T,p}$ of
\begin{equation}\label{rb:eq:coupling}
 \frac{(uX-vY)^z(vX+uY)^{T-z}}{\sqrt{z!(T-z)!}}.
\end{equation}
Descendants in every copy are obtained by normalized lowering, so the matching of bases between equal-spin copies is intertwining.
\end{lemma}
\begin{proof}
At $T=z$, the raising equation gives the adjacent ratio
\[
 -\sqrt{\frac{(z-p)(2j_2-z+p+1)}{(p+1)(2j_1-p)}}.
\]
With sign $(-1)^z$ at $p=0$, normalization yields the coefficients of $(uX-vY)^z/\sqrt{z!}$ in the limit. The change of variables $(X,Y)\mapsto(uX-vY,vX+uY)$ is orthogonal, so these polynomials have norm one in the normalized-monomial inner product. Normalized lowering from $T$ to $T+1$ divides the sum of the two lowering operators by
\[
 \sqrt{(T-z+1)\{2(j_1+j_2-z)-(T-z)\}}.
\]
In the limit this acts by multiplication by $(vX+uY)/\sqrt{T-z+1}$. Induction proves the statement for fixed descendants.
\end{proof}

\subsection{Rotational averages of explicit squares}
\label{rb:sec:squares}
We construct a positive rotational average and compare its three- and four-body terms with those of $H_q^2$. This section specifies the finite inequalities needed; Subsection~\ref{rb:app:certificate} proves them by exact rational arithmetic.

Haar measure $dg$ on $SU(2)$ is normalized to have total mass one. For an operator $A$, write $A(g)=U(g)AU(g)^\dagger$, with the appropriate Fock representation. Schur averaging sends a rank-one operator on an irreducible spin-$j$ space to its trace times $I/(2j+1)$. For several equivalent copies, it acts in the same way on the spin factor and retains the matrix between copies.

For each of the seven rows in Subsection~\ref{rb:app:certificate}, the integer
$t_\rho$ lies in $\{0,\ldots,7\}$. Set
\[
 \mathcal S_\rho=\{(p,j)\in\mathbb Z^2:0\le p\le7,\ 0\le j\le8,\
                         0\le p+j-t_\rho\le8\}.
\]
The row specifies real coefficients $\lambda_\rho$ and
$\alpha^\rho_{pj}$ for $(p,j)\in\mathcal S_\rho$; unlisted entries are zero.
Every coefficient sum for row $\rho$ below is over $\mathcal S_\rho$.
Our comparison operator is
\begin{equation}\label{rb:eq:F}
 K=\sum_{\rho=1}^7F_\rho^\dagger F_\rho,\qquad
 F_\rho=\lambda_\rho B_{t_\rho}
       +\sum_{(p,j)\in\mathcal S_\rho}\alpha^{\rho}_{pj}\,
                                  c_{p+j-t_\rho}^\dagger c_jB_p.
\end{equation}
Thus only modes $0,\ldots,8$ occur in these auxiliary squares.

For a single row, omit $\rho$ and put $i=p+j-t$, $k=p'+l-t$. Define
\[
 O_{p,j,k}=B(v_p\wedge e_j\wedge e_k)=c_kc_jB_p.
\]
Reordering $c_ic_k^\dagger=\delta_{ik}-c_k^\dagger c_i$ gives the two-, three-, and four-body terms of $F^\dagger F$:
\begin{align}
 &\lambda^2 B_t^\dagger B_t,\label{rb:eq:normal2}\\
 &\sum_{p,j}\lambda\alpha_{pj}\left[(c_iB_t)^\dagger c_jB_p
                         +(c_jB_p)^\dagger c_iB_t\right]
   +\sum_{p,j;p',l}\alpha_{pj}\alpha_{p'l}\delta_{ik}
                         (c_jB_p)^\dagger c_lB_{p'},\label{rb:eq:normal3}\\
 &-\sum_{p,j;p',l}\alpha_{pj}\alpha_{p'l}
                         O_{p,j,k}^\dagger O_{p',l,i}.\label{rb:eq:normal4}
\end{align}
Let $A_2,A_3,A_4$ denote these terms, summed over rows and averaged with factor $d_a$. Thus
\begin{equation}\label{rb:eq:Adecomp}
 \begin{aligned}
 \mathcal K_q&:=d_a\int K(g)\,dg=A_2+A_3+A_4,\qquad A_2=\eta H,\\
 \eta&=\sum_\rho\lambda_\rho^2=\frac{93527408868499}{10^{14}}.
 \end{aligned}
\end{equation}

\subsubsection{The three-body comparison}
Before wedging, \eqref{rb:eq:normal3} is a matrix on $V_q\otimes U_q$. Its nonzero entries preserve $T=p+j$, and $T\le15$. Averaging therefore gives
\begin{equation}\label{rb:eq:A3}
 A_3=\mathcal L_3\left(W_3\sum_{z=0}^{15}\frac{d_a}{d_{3,z}}e_z^q R_{3,z}W_3^\dagger\right),
\end{equation}
for all sufficiently large $q$, where $e_z^q$ is the trace of the unaveraged matrix in that spin block. By Lemma~\ref{rb:lem:coupling}, with $u=1/\sqrt3$, these traces converge to
\begin{equation}\label{rb:eq:ez}
 e_z=\sum_\rho\sum_{T=\max(z,t_\rho)}^{15}
 \left[2\lambda_\rho s_{z,T,t_\rho}
                   \sum_{p+j=T}\alpha^{\rho}_{pj}s_{z,T,p}
       +\left(\sum_{p+j=T}\alpha^{\rho}_{pj}s_{z,T,p}\right)^2\right],
\end{equation}
where $s=s^{(1/\sqrt3)}$ and empty inner sums are zero. The condition $i=k$ in \eqref{rb:eq:normal3} is exactly the equality of the two values of $T$.

The exact calculation in Subsection~\ref{rb:app:certificate} proves
\begin{equation}\label{rb:eq:3certificate}
 e_z<\frac32(3z-1)(-1/2)^z,
 \qquad z\in\{2,4,5,\ldots,15\}.
\end{equation}
There are finitely many strict inequalities, $d_{3,z}/d_a\to3/2$, and the omitted blocks $z=0,1,3$ are annihilated by $W_3$. Consequently, for all sufficiently large $q$,
\begin{equation}\label{rb:eq:A3bound}
 A_3\le\mathcal L_3\left(W_3\sum_{z=0}^{15}q_z(q)R_{3,z}W_3^\dagger\right).
\end{equation}

\subsubsection{The four-body comparison before taking limits}
Let
\[
 \widehat W:V_q\otimes\bigwedge^2U_q\longrightarrow\bigwedge^4U_q
\]
be the wedge map. The second tensor factor decomposes into spins $q-r$, with $r$ odd, $1\le r\le q$. In its lowering bases use the sign of Lemma~\ref{rb:lem:coupling}; at $r=1$ this is precisely the basis $v_p$.

A copy of spin $2q-1-D$ is obtained by coupling $a=q-1$ with $q-r$ at deficit $D-r$. Denote its vector at total deficit $T$ by $\rbket{D,T;r}_q$. Its physical orbital-index sum is $1+T$, and it is a highest vector when $T=D$. For fixed $D,T$ and sufficiently large $q$, the allowed copy labels are
\[
 \mathcal R_D=\{1,3,5,\ldots\}\cap[1,D],\qquad D\le T.
\]
Define the real finite-flux coefficients by
\[
 \rbket{D,T;r}_q=\sum_{\substack{p+j+k=T\\j<k}}
 S_r^q(D,T;p,j,k)\,v_p\otimes(e_j\wedge e_k),
\]
and extend them antisymmetrically in $j,k$. By two applications of Lemma~\ref{rb:lem:coupling}, they converge to
\begin{equation}\label{rb:eq:S}
 S_r(D,T;p,j,k)=\sqrt2\,
 s^{(1/\sqrt2)}_{r,j+k,j}\,
 s^{(1/\sqrt2)}_{D-r,T-r,p},
\end{equation}
interpreted as zero if $r>j+k$. The formula holds for either ordering of $j,k$, by antisymmetry. It follows by applying Lemma~\ref{rb:lem:coupling} first to the two single-particle spins and then to the two pair spins.

After averaging \eqref{rb:eq:normal4}, its spin-$2q-1-D$ block before wedging is a copy matrix $E_D^q$ times the spin identity, multiplied by $d_a/d_D$, where
\[
 d_D=4q-1-2D.
\]
Only $1\le D\le23$ contribute: every unaveraged term has
\[
 T=p+j+k=p'+l+i\le23.
\]
For these finitely many $D$, entrywise convergence gives $E_D^q\to E_D$, with
\begin{equation}\label{rb:eq:E}
 (E_D)_{rs}=-\sum_\rho\sum_{\substack{p,j;p',l\\T\ge D}}
 \alpha^{\rho}_{pj}\alpha^{\rho}_{p'l}
 S_r(D,T;p,j,k)S_s(D,T;p',l,i),
\end{equation}
where $i=p+j-t_\rho$, $k=p'+l-t_\rho$, and $T=p+j+k$. The exact finite-flux matrix $E_D^q$ is given by the same formula with each $S_r$ replaced by $S_r^q$.

The target term $\mathcal L_4(W_4W_4^\dagger)$ is obtained from the $r=1$ summand, since that summand is $V_q\otimes V_q$. Thus its copy matrix, in the units just used, is $(d_D/d_a)\chi$, where
\[
 \chi_{rs}=\mathbf1_{r=s=1},\qquad d_D/d_a\longrightarrow2.
\]
The summands of the target with $1\le D\le23$ equal
$\mathcal L_4(W_4R_q^{\mathrm{lo}}W_4^\dagger)$.
Since $A_4$ has no other spin blocks, the comparison with $A_4$ uses only
this part of the target. We keep
$\mathcal L_4(W_4R_q^{\mathrm{hi}}W_4^\dagger)$ separately.

Set
\begin{align}
 w_{Dr}^q&=\widehat W\rbket{D,D;r}_q,\label{rb:eq:wq}\\
 w_{Dr}^*&=\sum_{\substack{p+j+k=D\\j<k}}
 S_r(D,D;p,j,k)\,v_p^*\wedge e_j\wedge e_k,\qquad
 (G_D)_{rs}=\rbip{w_{Dr}^*}{w_{Ds}^*}.\label{rb:eq:wstar}
\end{align}
The star is a limit label, not an adjoint. The exact certificate is
\begin{equation}\label{rb:eq:4certificate}
 G_D(2\chi-E_D+\varepsilon I)G_D\ge0,
 \qquad 1\le D\le23,\qquad \varepsilon=3\cdot10^{-6}.
\end{equation}
We next show how these finite inequalities give the stated bound uniformly
in particle number. Subsection~\ref{rb:app:certificate} supplies their
integer data and exact rational verification.

\subsection{A relative error bound for finite spheres}
\label{rb:sec:transfer}
The limiting Gram matrices have kernels, so their positivity cannot be
transferred by assuming that their ranges vary continuously.  Instead, we
compare the annihilators themselves.  An exact diagonal change of variables
reduces every error to a fixed finite family of pair-generated annihilators.
This gives an error relative to $H_q$, uniformly on the full Fock space.

For the local argument, let $\mathcal F_{\mathrm{loc}}$ be the exterior
algebra on the 25 orthonormal modes $e_0,\ldots,e_{24}$.
For $q\ge24$ we identify it with the Fock space of the first 25 orbital
modes in $U_q$. By \eqref{rb:eq:pairstar}, each limiting pair vector $v_p^*$
with $p\le23$ lies in $\bigwedge^2\mathbb C^{25}$.
Contraction by these vectors is defined on $\mathcal F_{\mathrm{loc}}$
by the same exterior-algebra convention as before.

\begin{lemma}[Compression and annihilator bounds]\label{rb:lem:compression}
Let $W$ be a linear map between finite-dimensional Hilbert spaces,
$G=W^\dagger W$, and let $C$ be self-adjoint on the domain of $W$. Then
\[
 GCG\ge0\quad\Longleftrightarrow\quad WCW^\dagger\ge0.
\]
On $\mathcal F_{\mathrm{loc}}$, for any fixed finite collection
$A_b=c_kc_jB(v_p^*)$ with $0\le p\le23$ and $0\le j,k\le24$,
\begin{equation}\label{rb:eq:annihilatorbound}
 \sum_b\rbnorm{A_b\psi}^2\le C_0\rbip\psi{h_*\psi},
 \qquad h_*:=\sum_{p=0}^{23}B(v_p^*)^\dagger B(v_p^*),
\end{equation}
where $C_0$ depends only on the collection.
\end{lemma}
\begin{proof}
The two positivity statements both assert that the quadratic form of $C$
is nonnegative on $\operatorname{ran} G=\operatorname{ran} W^\dagger$.  For the second assertion, use
$\rbnorm{c_kc_jB(v_p^*)\psi}\le\rbnorm{B(v_p^*)\psi}$ and sum; one may take
$C_0$ to be the largest multiplicity of a pair label $p$.
\end{proof}

\begin{lemma}[Uniform transfer]\label{rb:lem:transfer}
For $q\ge24$ and $1\le D\le23$, put
\[
 C_D^q=(d_D/d_a)\chi-E_D^q+\varepsilon I,
 \qquad h_q=\sum_{p=0}^{23}B_p^\dagger B_p.
\]
Assuming \eqref{rb:eq:4certificate}, there is a scalar $\rho_q\ge0$,
independent of $D$ and tending to zero as $q\to\infty$, such that
\begin{equation}\label{rb:eq:relative}
 \sum_{r,s\in\mathcal R_D}(C_D^q)_{rs}
 B(w_{Dr}^q)^\dagger B(w_{Ds}^q)
 \ge-\rho_qh_q
\end{equation}
on the entire Fock space $\mathcal F_q$.
\end{lemma}
\begin{proof}
\emph{An exact change of variables.}
Take $q\ge24$.  All annihilators in \eqref{rb:eq:relative}, including those
in $h_q$, involve only the modes $0,\ldots,24$. First work on
$\mathcal F_{\mathrm{loc}}$. Define the positive diagonal operator
\[
 \Delta_q\rbket A=\left(\prod_{x\in A}d_q(x)\right)\rbket A,
 \qquad d_q(x)=\sqrt{(q)_x/q^x},
\]
on each basis wedge $\rbket A$.  This local Fock space is fixed as $q$
varies.  Thus $\Delta_q$ and $\Delta_q^{-1}$ converge in norm to $I$;
also $0<\Delta_q\le I$.

The pair coefficient formula \eqref{rb:eq:pair} gives
\[
 v_p(x,y)=r_p(q)d_q(x)d_q(y)v_p^*(x,y),
 \qquad r_p(q)=\sqrt{(2q)^p/(2q-2)_p}\ge1.
\]
Checking the factors attached to the removed orbital labels on a basis
wedge gives the exact identities
\begin{align}
 B_p&=r_p(q)\Delta_q^{-1}B(v_p^*)\Delta_q,
       \label{rb:eq:pairconjugation}\\
 B(v_p\wedge e_j\wedge e_k)
 &=\frac{r_p(q)}{d_q(j)d_q(k)}\,
   \Delta_q^{-1}B(v_p^*\wedge e_j\wedge e_k)\Delta_q.
       \label{rb:eq:fourconjugation}
\end{align}
These are identities on the local Fock space, with the inverse diagonal
acting on the output particle sector.  Repeated removed labels make both
sides zero.  In particular, solving \eqref{rb:eq:pairconjugation} for
$B(v_p^*)\Delta_q$ and using $r_p(q)\ge1$ and $\rbnorm{\Delta_q}\le1$
gives
\begin{equation}\label{rb:eq:pairenergytransfer}
 \Delta_q h_*\Delta_q\le h_q.
\end{equation}

\emph{The error is controlled by pair energy.}
Fix $D$.  Index a finite family by
$b=(p,j,k)$ with $p+j+k=D$ and $j<k$, and put
\[
 A_b=B(v_p^*\wedge e_j\wedge e_k),\qquad
 \beta_{rb}^q=S_r^q(D,D;p,j,k)
                \frac{r_p(q)}{d_q(j)d_q(k)}.
\]
Equations \eqref{rb:eq:wq} and \eqref{rb:eq:fourconjugation} yield
\[
 B(w_{Dr}^q)=\Delta_q^{-1}
             \left(\sum_b\beta_{rb}^qA_b\right)\Delta_q.
\]
The coefficients are real, and Lemma~\ref{rb:lem:coupling} gives
$\beta_{rb}^q\to\beta_{rb}^*:=S_r(D,D;p,j,k)$.  Hence the finite
real symmetric matrices
\[
 M_q=(\beta^q)^TC_D^q\beta^q
 \quad\hbox{converge to}\quad
 M_*=(\beta^*)^T(2\chi-E_D+\varepsilon I)\beta^*.
\]
The quadratic expression in \eqref{rb:eq:relative} is consequently
\[
 \Delta_q\left(\sum_{b,c}(M_q)_{bc}
                 A_b^\dagger\Delta_q^{-2}A_c\right)\Delta_q.
\]
Its limiting expression before the outer diagonals is
\begin{equation}\label{rb:eq:limitpositive}
 \sum_{b,c}(M_*)_{bc}A_b^\dagger A_c
 =\sum_{r,s}(2\chi-E_D+\varepsilon I)_{rs}
       B(w_{Dr}^*)^\dagger B(w_{Ds}^*)\ge0.
\end{equation}
Indeed, the certificate \eqref{rb:eq:4certificate} and
Lemma~\ref{rb:lem:compression}, applied to the map with columns $w_{Dr}^*$,
give positivity before lifting; the lift $\mathcal L_4$ preserves positivity.

For clarity, the perturbation estimate does not require $M_*\ge0$.
Let $R_q$ be the block matrix with entries
\[
 (R_q)_{bc}=(M_q)_{bc}\Delta_q^{-2}-(M_*)_{bc}I.
\]
Its norm tends to zero, since both its index set and the local Fock space
are fixed.  For every vector $\phi$,
\[
 \left|\sum_{b,c}\rbip{A_b\phi}{(R_q)_{bc}A_c\phi}\right|
 \le\rbnorm{R_q}\sum_b\rbnorm{A_b\phi}^2
 \le C_0\rbnorm{R_q}\rbip\phi{h_*\phi}.
\]
Combine this estimate with \eqref{rb:eq:limitpositive}, take
$\phi=\Delta_q\psi$, and apply \eqref{rb:eq:pairenergytransfer}.  This
proves \eqref{rb:eq:relative} on the local Fock space with a coefficient
tending to zero.  Taking a maximum over the fixed set $1\le D\le23$
gives a common $\rho_q$.

\emph{Additional particles are spectators.}
Order the 25 local modes before all other modes in the Fock
factorization.  Every operator in \eqref{rb:eq:relative} acts on the local
factor alone.  Tensoring the local inequality with the identity proves
it on $\mathcal F_q$, with the same $\rho_q$, regardless of the number of
occupied modes outside this factor.
\end{proof}

The allowance $\varepsilon I$ in the certificate must also be paid for in pair
energy.  Here it is useful to sum over the orthonormal highest vectors
\emph{before} estimating their wedges.  Bessel's inequality then avoids
an extra factor equal to the number of spin copies.

\begin{proposition}[Four-body estimate]\label{rb:prop:A4}
For the coefficients in Subsection~\ref{rb:app:certificate},
\begin{equation}\label{rb:eq:A4bound}
 A_4\le\mathcal L_4(W_4R_q^{\mathrm{lo}}W_4^\dagger)+(1222\varepsilon+o(1))H,
 \qquad q\to\infty,
\end{equation}
uniformly on $\mathcal F_q$.
\end{proposition}
\begin{proof}
Haar averaging the matrix unit
$\rbket{D,D;r}_q\rbbra{D,D;s}_q$ gives the corresponding copy matrix unit
tensored with the spin identity, divided by $d_D$.  Therefore, on the
blocks $1\le D\le23$, the low-block target minus $A_4$, with the $\varepsilon I$ allowance
added, is
\[
 d_a\sum_{D=1}^{23}\int\sum_{r,s}(C_D^q)_{rs}
 B(w_{Dr}^q(g))^\dagger B(w_{Ds}^q(g))\,dg.
\]
Lemma~\ref{rb:lem:transfer}, conjugated by each rotation, bounds this below
by $-o(1)H$, because Schur averaging gives
\[
 d_a\int h_q(g)\,dg=24H.
\]
The number of blocks in the comparison is fixed, so this error remains uniform
in particle number.

To bound the allowance, fix $D$ and write
\[
 \rbket{D,D;r}_q
   =\sum_{\substack{p+j+k=D\\j<k}}S_{rb}^q\,
                     v_p\otimes(e_j\wedge e_k),
 \qquad b=(p,j,k).
\]
The uncoupled vectors on the right are orthonormal.  So are the highest
vectors on the left: distinct $r$ belong to orthogonal spin summands of
the second pair factor.  Thus $S^q(S^q)^T=I$.  Apply the contraction
$S^q\otimes I$ to the Hilbert-space-valued vector
$(c_kc_jB_p\psi)_b$.  It follows that
\[
 \sum_{r\in\mathcal R_D}\rbnorm{B(w_{Dr}^q)\psi}^2
 \le\sum_{\substack{p+j+k=D\\j<k}}\rbnorm{c_kc_jB_p\psi}^2
 \le\sum_{\substack{p+j+k=D\\j<k}}\rbnorm{B_p\psi}^2.
\]
This uses orthogonality before wedging; the vectors $w_{Dr}^q$ themselves
need not be orthogonal.  The number of triples on the right is
\[
 n_D=\sum_{n=1}^D\left\lceil\frac n2\right\rceil
     =\left\lfloor\frac{(D+1)^2}{4}\right\rfloor.
\]
Conjugate by rotations and integrate.  Each pair norm contributes
$\rbip\psi{H\psi}/d_a$, so the total allowance is bounded by
\[
 \varepsilon\sum_{D=1}^{23}n_DH
 =\varepsilon\left(\sum_{m=1}^{12}m^2+
                 \sum_{m=1}^{11}m(m+1)\right)H
 =1222\varepsilon H.
\]
Removing the allowance proves \eqref{rb:eq:A4bound}. No block of
$R_q^{\mathrm{hi}}$ has entered this estimate.
\end{proof}

\subsection{The three-body tail and the retained-block inequality}
\label{rb:sec:finish}
The preceding comparison has used only the low four-body target.
It remains to control the three-body blocks above $z=15$ by pair energy
and combine the estimates.

\begin{lemma}[Three-body tail]\label{rb:lem:tail}
For sufficiently large $q$,
\begin{equation}\label{rb:eq:tailbound}
 \mathcal L_3\left(W_3\sum_{z=16}^q q_z(q)R_{3,z}W_3^\dagger\right)
 \ge-\delta_qH,
\end{equation}
where
\begin{equation}\label{rb:eq:delta}
 \delta_q=\sum_{\substack{17\le z\le q\\z\text{ odd}}}
       \frac{d_{3,z}}{d_a}|q_z(q)|(z+1),\qquad
 \delta_q\longrightarrow\frac{61}{4096}.
\end{equation}
\end{lemma}
\begin{proof}
For $16\le z\le q$, the bracket in \eqref{rb:eq:q} is positive, since it is at least $2z-2$. Thus only odd $z$ contribute negatively. A normalized highest vector of $R_{3,z}$ is $\sum_{p=0}^zs_pv_p\otimes e_{z-p}$. Its orbit resolves that spin projector with factor $d_{3,z}$. After wedging and lifting, its quadratic form is
\[
 \rbnorm{\sum_{p=0}^zs_pc_{z-p}(g)B_p(g)\psi}^2
 \le\sum_{p=0}^z\rbnorm{B_p(g)\psi}^2.
\]
Averaging each summand proves \eqref{rb:eq:tailbound} with \eqref{rb:eq:delta}.

For $q\ge4$, the ratios in $(q)_z/(2q-2)_z$ are nonincreasing, so
\[
 \frac{(q)_z}{(2q-2)_z}\le\left(\frac q{2q-2}\right)^z
 \le(2/3)^z,\qquad \frac{d_{3,z}}{d_a}\le2.
\]
The summands, extended by zero for $z>q$, are therefore dominated by
$6z(z+1)(2/3)^z$. Dominated convergence and \eqref{rb:eq:qlimit} give
\[
 \lim_{q\to\infty}\delta_q
 =\frac32\sum_{z=17,19,\ldots}(3z-1)(z+1)2^{-z}
 =\frac{61}{4096}.
\]
The last identity follows by writing $z=17+2m$ and using the geometric-series identities for $\sum t^m$, $\sum mt^m$, and $\sum m^2t^m$ at $t=1/4$.
\end{proof}



\begin{proof}[Proof of Proposition~\ref{prop:retained-blocks}]
Write $T_{\mathrm{lo}}=\mathcal L_4(W_4R_q^{\mathrm{lo}}W_4^\dagger)$
and $T_{\mathrm{hi}}=\mathcal L_4(W_4R_q^{\mathrm{hi}}W_4^\dagger)$.
The normal-square identity, the three-body comparison
\eqref{rb:eq:A3bound}, and the tail estimate give
\[
 H^2\ge (1-\delta_q)H+A_3+T_{\mathrm{lo}}+T_{\mathrm{hi}}.
\]
By \eqref{rb:eq:A4bound},
$T_{\mathrm{lo}}\ge A_4-(1222\varepsilon+o(1))H$.
Using $\mathcal K_q=\eta H+A_3+A_4\ge0$ from
\eqref{rb:eq:Adecomp}, we obtain
\[
 H^2\ge\mathcal K_q+
 (1-\eta-\delta_q-1222\varepsilon-o(1))H+T_{\mathrm{hi}}.
\]
The coefficient of $H$ tends to
\begin{equation}\label{rb:eq:finalmargin}
 1-\frac{93527408868499}{10^{14}}-\frac{61}{4096}
   -1222\frac3{10^6}
 =\frac{4616733319001}{10^{14}}>\frac1{25}.
\end{equation}
For fixed $0<c_0<1/25$, choose $q$ large enough that this coefficient
exceeds $c_0$ and discard the positive operator $\mathcal K_q$.
This is \eqref{rb:eq:retained-blocks}. The matrix comparisons involved
finitely many fixed orbital modes, the transfer estimates held on full
Fock space, and the tail estimate was uniform in particle number.
Therefore the same threshold applies simultaneously to every sector.
\end{proof}


\subsection{The finite rational certificate}
\label{rb:app:certificate}
We finish by verifying the two finite inequalities
\eqref{rb:eq:3certificate} and \eqref{rb:eq:4certificate} used above.
The seven integer rows and all rational formulas are given here.
Starting from the integer
rows in Subsection~\ref{rb:app:rows}, evaluate the rational three-body traces in
Subsection~\ref{rb:app:three-arithmetic} and the rational four-body matrices in
Subsection~\ref{rb:app:four-arithmetic}. The margin and coefficient-sign tables in
those sections record the resulting checks. The formulas below determine
every value using only integers and rational numbers.

\subsubsection{Row data}
\label{rb:app:rows}
Let $P=10^7$. The first two numbers of each row are $t$ and $\ell=P\lambda$. An entry $pj:a$ specifies the integer $a_{pj}$, with $p$ and $j$ parsed as separate single-digit labels; omitted entries are zero. Continuation lines belong to the preceding row. Define
\begin{equation}\label{rb:eq:alphadata}
 \alpha_{pj}=\frac{a_{pj}}P
       \sqrt{\frac{2^p\binom{p+j}{p}}{2^t\binom{p+j}{t}}}
 =\frac{a_{pj}}P\sqrt{\frac{2^{p-t}t!(p+j-t)!}{p!j!}}.
\end{equation}
Every nonzero entry has $0\le p+j-t\le8$.
\par\noindent\begin{minipage}{\linewidth}
\lstinputlisting[basicstyle=\ttfamily\scriptsize,columns=fullflexible,keepspaces=true,breaklines=true,frame=single]{../verification/rows.txt}
\end{minipage}\par
\smallskip
\noindent In particular $P^2\eta=\sum_\rho\ell_\rho^2=93527408868499$.

\subsubsection{Three-body arithmetic}
\label{rb:app:three-arithmetic}
For $c\in\{1,2\}$, an integer $T\ge0$, and integers $z,p$, define
\begin{equation}\label{rb:eq:U}
 \mathsf U(c,z,T,p)=
 \sum_{h=\max(0,p+z-T)}^{\min(p,z)}
 (-1)^{z-h}\binom zh\binom{T-z}{p-h}c^{p-h},
\end{equation}
and set it to zero if $z$ or $p$ lies outside $[0,T]$. Expanding \eqref{rb:eq:coupling} gives
\[
 s^{(1/\sqrt3)}_{z,T,p}=\mathsf U(2,z,T,p)
 \sqrt{\frac{2^z\binom Tz}{3^T2^p\binom Tp}}.
\]
Substitution in \eqref{rb:eq:ez} and \eqref{rb:eq:alphadata} cancels all radicals:
\begin{equation}\label{rb:eq:erational}
 P^2e_z=\sum_\rho\sum_{T=\max(z,t_\rho)}^{15}
 \frac{2^z\binom Tz}{3^T2^{t_\rho}\binom T{t_\rho}}
 X_{\rho,T}\left(X_{\rho,T}+2\ell_\rho\mathsf U(2,z,T,t_\rho)\right),
\end{equation}
where
\[
 X_{\rho,T}=\sum_{p+j=T}a^\rho_{pj}\mathsf U(2,z,T,p),
\]
with value zero on empty levels. For each $z\in\{2,4,5,\ldots,15\}$,
evaluate \eqref{rb:eq:erational} and form the exact rational margin
$m_z=\frac32(3z-1)(-1/2)^z-e_z$. Write $10^6m_z=n_z/d_z$ with integers $n_z,d_z$ and $d_z>0$. Integer division gives the floors $\lfloor n_z/d_z\rfloor$ listed below:
\begin{center}
\small
\setlength{\tabcolsep}{4pt}
\begin{tabular}{rrrrrrrrrrrrrr}
\toprule
$z$&2&4&5&6&7&8&9&10&11&12&13&14&15\\
\midrule
floor&1298313&211317&249&250&249&250&250&250&249&250&250&6470&250\\
\bottomrule
\end{tabular}
\end{center}
All are positive, proving \eqref{rb:eq:3certificate}.

\subsubsection{Four-body arithmetic}
\label{rb:app:four-arithmetic}
Fix $1\le D\le23$; all copy indices below lie in $\mathcal R_D$. Put
\[
 u_r=\frac1{r!(D-r)!},\qquad
 (E_D)_{rs}=\sqrt{u_ru_s}\,Y_{rs},\qquad
 (G_D)_{rs}=\sqrt{u_ru_s}\,Z_{rs}.
\]
The matrices $Y,Z$ are rational. Define
\[
 \mathsf V_r(D,T;p,j,k)=
 \mathsf U(1,r,j+k,j)\mathsf U(1,D-r,T-r,p),
\]
with value zero for $r>j+k$. Two applications of \eqref{rb:eq:coupling} give
\begin{equation}\label{rb:eq:Srational}
 S_r(D,T;p,j,k)=\mathsf V_r(D,T;p,j,k)
 \sqrt{2^{1-(j+k)-T+r}\frac{j!k!p!\,u_r}{(T-D)!}}.
\end{equation}
Therefore
\begin{equation}\label{rb:eq:Yrational}
 \begin{split}
 P^2Y_{rs}=-\sum_\rho\sum_{\substack{p,j;p',l\\T\ge D}}
 &a^\rho_{pj}a^\rho_{p'l}\,
 \mathsf V_r(D,T;p,j,k)\mathsf V_s(D,T;p',l,i)\\
 &\times2^{1-2T+p+p'-t_\rho+(r+s)/2}
       \frac{i!k!t_\rho!}{(T-D)!},
 \end{split}
\end{equation}
where $i=p+j-t_\rho$, $k=p'+l-t_\rho$, and $T=p+j+k$. Both entry sums are ordered. To verify the cancellation, multiply the two factors in \eqref{rb:eq:Srational} by the two factors \eqref{rb:eq:alphadata}; their remaining square root is exactly the factorial and power-of-two factor displayed in \eqref{rb:eq:Yrational}.

For an increasing quadruple $A$ of nonnegative indices with sum $D+1$, put $f(A)=\prod_{b\in A}b!$. If $(x,y,j,k)$ lists its elements, let $\sigma(x,y,j,k)$ be the sign of sorting this list into increasing order. In each summand below put $p=x+y-1$, and define
\begin{equation}\label{rb:eq:Lrational}
 L_r(A)=2^{(1-2D+r)/2}
 \sum_{\substack{x<y\text{ in }A\\\{j,k\}=A\setminus\{x,y\},\ j<k}}
 \sigma(x,y,j,k)(x-y)\,p!j!k!\,
 \mathsf V_r(D,D;p,j,k).
\end{equation}
Combining \eqref{rb:eq:pairstar} and \eqref{rb:eq:Srational}, the coefficient of $w_{Dr}^*$ on the Slater wedge $A$ is $\sqrt{u_r/f(A)}L_r(A)$. Thus
\begin{equation}\label{rb:eq:Zrational}
 Z_{rs}=\sum_{\substack{A\text{ increasing quadruple}\\\sum_{b\in A}b=D+1}}
                   \frac{L_r(A)L_s(A)}{f(A)}.
\end{equation}
All exponents of two in these formulas are integers, since $r,s$ are odd.

Let $D_u=\operatorname{diag}(u_r)$ and $D_{\sqrt u}=\operatorname{diag}(\sqrt{u_r})$. The matrix in \eqref{rb:eq:4certificate} equals $D_{\sqrt u}MD_{\sqrt u}$, where
\begin{equation}\label{rb:eq:M}
 M=ZBZ,\qquad
 B=\operatorname{diag}\bigl((2\mathbf1_{r=1}+\varepsilon)u_r\bigr)-D_uYD_u.
\end{equation}
Thus it suffices to check $M\ge0$ using rational arithmetic.

To check positivity for each $D=1,\ldots,23$, form $Y$ and $Z$ from
\eqref{rb:eq:Yrational} and \eqref{rb:eq:Zrational}, then form $B$ and $M$ from
\eqref{rb:eq:M}, using $\varepsilon=3/10^6$. All sums range over the printed rows and
the finite index sets specified above. Compute with exact fractions
throughout. If $d=|\mathcal R_D|$, initialize $J=I_d$. For
$b=1,\ldots,d$, compute in order
\begin{equation}\label{rb:eq:recurrence}
 X=MJ,\qquad c_b=\frac{\operatorname{tr} X}{b},\qquad J=c_bI_d-X.
\end{equation}
Then $\det(xI+M)=x^d+c_1x^{d-1}+\cdots+c_d$. For example, this recurrence follows by multiplying the polynomial adjugate of $xI+M$ by $xI+M$ and comparing coefficients, together with the derivative-of-determinant trace identity. Exact evaluation of \eqref{rb:eq:U}, \eqref{rb:eq:Yrational}, \eqref{rb:eq:Lrational}, \eqref{rb:eq:Zrational}, and \eqref{rb:eq:recurrence} gives:
\begin{center}
\begin{tabular}{rl@{\qquad}rl@{\qquad}rl}
\toprule
$D$&signs&$D$&signs&$D$&signs\\
\midrule
1&\texttt{0}&9&\texttt{++000}&17&\texttt{++++++000}\\
2&\texttt{0}&10&\texttt{+0000}&18&\texttt{+++++0000}\\
3&\texttt{00}&11&\texttt{+++000}&19&\texttt{+++++++000}\\
4&\texttt{00}&12&\texttt{++0000}&20&\texttt{++++++0000}\\
5&\texttt{+00}&13&\texttt{++++000}&21&\texttt{++++++++000}\\
6&\texttt{000}&14&\texttt{+++0000}&22&\texttt{+++++++0000}\\
7&\texttt{+000}&15&\texttt{+++++000}&23&\texttt{+++++++++000}\\
8&\texttt{+000}&16&\texttt{++++0000}&&\\
\bottomrule
\end{tabular}
\end{center}
The signs list the coefficients after the leading $1$, in descending degree.
For each computed coefficient, choose a positive denominator and inspect
its integer numerator: \texttt{+} denotes a positive numerator and
\texttt{0} denotes zero. Every coefficient is nonnegative. Hence
\[
 \det(xI+M)>0\qquad(x>0).
\]
Since $M$ is real symmetric, a negative eigenvalue would give a positive
root of this polynomial. Therefore $M\ge0$, proving
\eqref{rb:eq:4certificate}.


\endgroup
